%% kapa: exam=1 round=22 year=2011 period=1 form=A subject=economics subject_ko=경제원론 seq=9 number=49
%% source: Q-Net 감정평가사 기출문제 (cst00309) artlSeq=5036790 file=extracted/1008/5036790/2011년 제22회 감정평가사 1차시험 기출문제/1교시.hwp
%% conversion: kordoc 4.24.51 -> tools/gen1.py (auto); review=proofread-sonnet; flags=none
\begin{question}{49}
기대효용이론에 관한 설명으로 옳은 것은? (단, U는 효용수준, M은 자산액)
\begin{choices}
  \item 폰 노이만-모겐스턴(Von Neumann-Morgenstern) 효용함수에서 효용은 서수적 의미만 갖는다.
  \item 甲이 가지고 있는 복권 상금의 기대가치는 500이고 이 복권을 최소 450에 팔 용의가 있다면, 50을 甲의 위험프리미엄(risk premium)으로 볼 수 있다.
  \item 위험기피자는 기대가치가 0인 복권을 구입할 것이다.
  \item 위험선호자는 기대가치가 0인 보험에 가입할 것이다.
  \item 乙의 폰 노이만-모겐스턴 효용함수가 U=M\textsuperscript{1.5}로 주어졌다면, 乙은 위험기피자이다.
\end{choices}
\end{question}
\answer{2}
